集合~$A$ 上的一个\emph{关系}~$R$，是指~$A$ 中元素之间的一种关联方式。
如果该关系在 $x$ 与~$y$ 之间\emph{成立}，我们就记作 $Rxy$。
形式上说，我们可以把 $R$ 看作满足~$Rxy$ 的所有有序对
$\tuple{x,y} \in A^2$ 的集合。
小于、大于、等于、整除、长度相同、子集关系、大小相同等等，
都是关系的重要例子（在数的集合、字符串的集合或集合的集合上）。
\emph{图}是关系的一种通用的可视化表示方式。但图也可以看作一个二元关系（即\emph{边}关系）
连同其底层的\emph{顶点}集合。

有些关系具有某些共同特征，这使得它们特别有趣或有用。
关系~$R$ 是\emph{自反的}，如果每个元素都与自身有~$R$ 关系；
是\emph{对称的}，如果对任意 $x$ 和~$y$，只要 $Rxy$ 成立，则 $Ryx$ 也成立；
是\emph{传递的}，如果 $Rxy$ 和 $Ryz$ 能保证~$Rxz$。
同时具备这三个性质的关系称为\emph{等价关系}。
关系~$R$ 是\emph{反对称的}，如果 $Rxy$ 和 $Ryx$ 能保证~$x=y$。
\emph{偏序}就是那些自反、反对称且传递的关系。
\emph{线性序}则是满足如下条件的偏序：对任意 $x$ 和~$y$，要么 $Rxy$，要么 $x=y$，要么 $Ryx$。
（一般来说，具有这种性质的关系称为\emph{连通的}。）

既然关系是（有序对的）集合，它们就可以像集合一样进行运算（例如，我们可以取关系的并和交）。
我们还可以将它们链式组合起来（即\emph{相对积}~$R \mid S$）。
如果将~$R$ 与自身作任意多次相对积，我们就得到~$R$ 的\emph{传递闭包}~$R^+$。